% Zone „Das kennst du schon“ – aus bank/umkehrfunktion/zone.jsonl (zusammenbau v0.1)
\einheitenkopf[zone][Kennst du schon]{Das kennst du schon}
\setcounter{aufgabe}{0}

%% TODO Titel: Ich-kann-Satz fehlt in der Bank (Platzhalter: Kettenname) – Nr. 1
%% TODO Anweisung: ein Satz über der ersten Teilaufgabe fehlt in der Bank – Nr. 1
\begin{aufgabe}{Eine Gleichung nach x auflösen und die Variablen tauschen, bei e-Funktionen mit ln (Exponentialgleichung)}
\begin{teile}
\teil Löse $y = 3x - 6$ nach x auf und berechne x für $y = 9$. x = \leerfeld
\teil Löse $y = -0{,}5x + 4$ nach x auf und berechne x für $y = 3$. x = \leerfeld
\end{teile}
\end{aufgabe}

%% TODO Titel: Ich-kann-Satz fehlt in der Bank (Platzhalter: Kettenname) – Nr. 2
%% TODO Anweisung: ein Satz über der ersten Teilaufgabe fehlt in der Bank – Nr. 2
\begin{aufgabe}{Wertebereich einer Funktion aus Monotonie und Grenzverhalten angeben, strenge Monotonie am Graphen oder über die Ableitung erkennen}
\begin{teile}
\teil $f(x) = 2x + 1$ auf dem Intervall $[0;\, 3]$. Gib den Wertebereich an.
\teil $f(x) = -x + 5$ auf dem Intervall $[-2;\, 3]$. Gib den Wertebereich an.
\end{teile}
\end{aufgabe}

%% TODO Titel: Ich-kann-Satz fehlt in der Bank (Platzhalter: Kettenname) – Nr. 3
%% TODO Anweisung: ein Satz über der ersten Teilaufgabe fehlt in der Bank – Nr. 3
\begin{aufgabe}{Punkte an einer Geraden spiegeln, Bildpunkte benennen, Spiegelachse als Symmetrieachse einer Figur}
\begin{teile}
\teil Spiegle den Punkt A(2 | 5) an der x-Achse. A' = \leerfeld
\teil Spiegle den Punkt C(1,5 | −2) an der Geraden $y = 1$. C' = \leerfeld
\end{teile}
\end{aufgabe}

%% TODO Titel: Ich-kann-Satz fehlt in der Bank (Platzhalter: Kettenname) – Nr. 4
%% TODO Anweisung: ein Satz über der ersten Teilaufgabe fehlt in der Bank – Nr. 4
\begin{aufgabe}{Integral einer Differenz als Fläche zwischen zwei Graphen deuten, Trapezformel und gleichschenklig-rechtwinkliges Dreieck}
\begin{teile}
\teil Ein Trapez hat die parallelen Seiten 6 cm und 4 cm und die Höhe 3 cm. Berechne den Flächeninhalt. \leerfeld
\teil Berechne den Inhalt der Fläche zwischen den Graphen von $y = x$ und $y = 0{,}5x^2$ von $x = 0$ bis $x = 2$. \leerfeld
\end{teile}
\end{aufgabe}

%% TODO Titel: Ich-kann-Satz fehlt in der Bank (Platzhalter: Kettenname) – Nr. 5
%% TODO Anweisung: ein Satz über der ersten Teilaufgabe fehlt in der Bank – Nr. 5
\begin{aufgabe}{Steigung einer Geraden und Steigung der gespiegelten Geraden (Kehrwert; Steigung minus eins senkrecht zu y = x)}
\begin{teile}
\teil Bestimme die Steigung der Geraden durch (0 | 1) und (2 | 7). m = \leerfeld
\teil Bestimme die Steigung der Geraden durch (−1 | 4) und (3 | 2). m = \leerfeld
\end{teile}
\end{aufgabe}

%% TODO Titel: Ich-kann-Satz fehlt in der Bank (Platzhalter: Kettenname) – Nr. 6
%% TODO Anweisung: ein Satz über der ersten Teilaufgabe fehlt in der Bank – Nr. 6
\begin{aufgabe}{Wertebereich einer Funktion aus Monotonie und Grenzverhalten angeben, strenge Monotonie am Graphen oder über die Ableitung erkennen – Fehler finden}
\begin{teile}
\teil Lena gibt für $f(x) = e^x + 1$ auf ganz ℝ den Wertebereich $[1;\, \infty[$ an. Finde den Fehler.
\end{teile}
\end{aufgabe}

%% TODO Titel: Ich-kann-Satz fehlt in der Bank (Platzhalter: Kettenname) – Nr. 7
%% TODO Anweisung: ein Satz über der ersten Teilaufgabe fehlt in der Bank – Nr. 7
\begin{aufgabe}{Wertebereich einer Funktion aus Monotonie und Grenzverhalten angeben, strenge Monotonie am Graphen oder über die Ableitung erkennen – selbst rechnen}
\begin{teile}
\teil $f(x) = 3 - e^{-x}$ auf dem Intervall $[0;\, \infty[$ ist streng steigend. Gib den Wertebereich an.
\end{teile}
\end{aufgabe}
